\documentclass[UTF8,fontset=none,11pt,a4paper]{ctexart}
\setCJKmainfont{Noto Serif CJK SC}
\setCJKsansfont{Noto Sans CJK SC}
\setCJKmonofont{Noto Sans CJK SC}
\usepackage[margin=22mm]{geometry}
\usepackage{amsmath,booktabs,graphicx,hyperref,float}
\hypersetup{colorlinks=true,linkcolor=blue,urlcolor=blue}
\graphicspath{{docs/reports/NCFV_convergence/}}
\makeatletter
\let\tableinput\@@input
\makeatother
\newcommand{\FineOrderLtwo}{2.434}
\newcommand{\FineOrderNominal}{2.418}
\newcommand{\FineErrorLtwo}{1.081036e-03}
\newcommand{\FineReduction}{47.70}
\newcommand{\FineInterpretation}{最细两级的实测收敛阶尚不能直接认定为稳定三阶；应以表中实测值为准。}

\title{高效格点有限体积格式：等熵涡网格收敛验证\\
       \large iv10、iv20、iv40、iv80，同一物理时间 $t=2$}
\author{DNDSR 独立 NCFV 模块}
\date{2026年9月7日}
\begin{document}
\maketitle

\section{结论与验证范围}
三套新增网格 iv10、iv20、iv80 均使用高效全微分积分模式计算至 $t=2$，
并与此前已完成的 iv40 高效模式计算作比较。四套网格使用同一求解器二进制，
本次没有修改任何空间或时间离散代码，也没有修改 DNDSR 的原有模块。
iv80 的密度 $L_2$ 误差为 \texttt{\FineErrorLtwo}，较 iv10 减小约
\FineReduction{} 倍。iv40 到 iv80 的实测 $L_2$ 收敛阶为
\FineOrderLtwo。

\FineInterpretation{}
这里报告的是本次网格族、此时刻和当前参数下的实际误差，
不要求与论文原程序的误差数值相等，也没有按目标阶数修正数据。
几何构造仍严格采用边中点、原始面顶点平均和原始体顶点平均。

\section{统一物理与数值设置}
\begin{center}
\begin{tabular}{ll}\toprule
计算域、周期长度 & $[0,10]\times[0,10]\times[0,4]$；$(10,10,4)$\\
初始涡心、强度 & $(5,5)$；$\beta=5$\\
背景状态 $(\rho,u,v,w,p)$ & $(1,1,1,0,1)$\\
比热比、黏性、限制器 & $\gamma=1.4$；黏性关闭，限制器关闭\\
高效积分模式 & \texttt{EfficientDifferential}\\
黎曼通量、时间推进 & \texttt{Roe\_M2}；SSPRK3\\
时间步 & 全域 CFL $=0.5$，$\Delta t\le0.01$，最后一步截短\\
物理终止时刻 & $t=2$，即完整周期 $10$ 的 $20\%$\\
\bottomrule\end{tabular}
\end{center}
六个边界分区全部使用平移周期条件。局部伪时间关闭，所有进程使用同一物理
步长。粗网格主要受 $0.01$ 上限限制；细网格的步长由 CFL 条件进一步缩小。
新计算并发执行，因此各算例墙钟时间不构成严格的并行加速比基准。

设 $\operatorname{wrap}_{10}$ 把坐标映射到最近周期像，解析解为
\[
X=\operatorname{wrap}_{10}(x-5-t),\quad
Y=\operatorname{wrap}_{10}(y-5-t),\quad
a=\frac{5}{2\pi}\exp\!\left(\frac{1-X^2-Y^2}{2}\right),
\]
\[
u=1-aY,\quad v=1+aX,\quad w=0,\quad
T=1-\frac{\gamma-1}{2\gamma}a^2,\qquad
\rho_{\mathrm{ex}}=T^{1/(\gamma-1)},\quad p_{\mathrm{ex}}=\rho_{\mathrm{ex}}^\gamma.
\]
后处理独立计算解析密度，并逐节点核对求解器导出的解析值。
该 Gaussian 涡采用周期延拓，尾部并非严格紧支撑；
当前误差验证不能据此推广为任意周期光滑解的阶数证明。

\section{网格与对偶控制体核验}
\begin{center}\small
\begin{tabular}{lrrrrrr}\toprule
网格 & 原始节点 & 三棱柱 & 周期自由度 & $h$ & 时间步数 & MPI\\\midrule
\tableinput{docs/reports/NCFV_convergence/mesh_rows.tex}
\bottomrule\end{tabular}
\end{center}
四套网格沿 $z$ 方向分别有 4、8、16、32 层，具有相同的计算域。
原始 CGNS 的旧边界描述转换为读取器支持的 PointRange/FaceCenter。
每套兼容副本的三组坐标及全部体/面连接数组均与转换前记录的 SHA-256 一致，
没有平滑、重划分或改变输入网格坐标。

对偶微体为每条 $(c,f,e,i)$ 拓扑链对应的四面体
\[
\operatorname{conv}(\boldsymbol x_i,\boldsymbol x_e,\boldsymbol x_f,\boldsymbol x_c),
\qquad
\boldsymbol x_e=\tfrac12(\boldsymbol x_i+\boldsymbol x_j),\
\boldsymbol x_f=\tfrac1{N_f}\sum_{j\in f}\boldsymbol x_j,\
\boldsymbol x_c=\tfrac1{N_c}\sum_{j\in c}\boldsymbol x_j .
\]
周期节点的各原始控制体份额合并成完整控制体。所有网格的份额体积和均为
$400$（浮点容差内），体、内部面和边界面的 Gauss 点存储计数均为零。
几何闭合误差及周期状态一致性见后表。

\section{误差口径与实测收敛阶}
自由度是对偶体平均守恒量，不能直接把它与解析节点点值比较。
本次使用求解器从对偶体均值恢复的节点密度 $\rho_i^{pt}$。
令 $V_i$ 为该原始节点在计算域内的对偶体份额体积：
\[
e_i=\rho_i^{pt}-\rho_{\mathrm{ex}}(\boldsymbol x_i,t),\qquad
L_1=\frac{\sum_iV_i|e_i|}{\sum_iV_i},\quad
L_2=\sqrt{\frac{\sum_iV_i e_i^2}{\sum_iV_i}},\quad
L_\infty=\max_i|e_i|.
\]
只读取每个原始节点唯一 owner 导出的 CSV，不重复计算 MPI ghost。
同一周期类各份额权重之和恰为完整周期控制体体积。后处理的三个误差范数
与求解器 MPI 汇总的结果分别核对到浮点舍入量级。

由于非结构网格自由度并不恰好每次增加八倍，主要采用实际自由度定义
\[
h_k=(400/N_{k,\mathrm{periodic}})^{1/3},\qquad
p_k=\frac{\log(E_{k-1}/E_k)}{\log(h_{k-1}/h_k)} .
\]
表中 $p$ 是上一行到本行的实测阶。若采用名义尺度 $h=10/n$、固定加密比
$2$，最细两级的 $L_2$ 阶为 \FineOrderNominal；
完整名义阶和面内尺度也保存在机器可读结果中。

\begin{center}\scriptsize
\begin{tabular}{lrrrrrr}\toprule
网格 & $L_1$ & $p_1$ & $L_2$ & $p_2$ & $L_\infty$ & $p_\infty$\\\midrule
\tableinput{docs/reports/NCFV_convergence/error_rows.tex}
\bottomrule\end{tabular}
\end{center}
\begin{figure}[H]\centering
\includegraphics[width=\linewidth]{convergence.pdf}
\caption{左：$t=2$ 密度误差；右：初始化及点值恢复误差。
灰线仅为斜率参照，不参与误差计算。}
\end{figure}
粗网格的涡心区域解析梯度变化较快，不能仅由粗网格对的阶数推断渐近阶数。
反之，也不能仅因算法设计目标为三阶，就把所有网格对的实测阶写为三阶。
本验证保留每一级的真实数据，并单独检查初始化和时间步影响。

\section{时间步与守恒检查}
另在 iv20 和 iv40 上将 CFL 与步长上限同时减半：
CFL $=0.25$、$\Delta t\le0.005$，其余物理和空间参数不变，同样计算至 $t=2$。
令 $D_2=\|\rho_{\Delta t}^{pt}-\rho_{\Delta t/2}^{pt}\|_{L_2(V)}$，
直接比较两个数值解，避免只比较误差范数时的抵消。
\begin{center}\scriptsize
\begin{tabular}{lrrrr}\toprule
网格 & 原步长 $L_2$ & 半步长 $L_2$ & $D_2$ & $100D_2/L_2$（\%）\\\midrule
\tableinput{docs/reports/NCFV_convergence/time_rows.tex}
\bottomrule\end{tabular}
\end{center}
这两套复算用于检查时间误差敏感性，不等于对 iv80 也完成了独立的时间阶数
验证。表中守恒漂移为 $|Q(2)-Q(0)|/\max(|Q(0)|,1)$，
周期差是五个保守变量在周期副本间差异的最大值。
\begin{center}\scriptsize
\begin{tabular}{lrrrrr}\toprule
网格 & 初始密度 $L_2$ & 质量漂移 & 总能量漂移 & 周期差 & 几何闭合\\\midrule
\tableinput{docs/reports/NCFV_convergence/check_rows.tex}
\bottomrule\end{tabular}
\end{center}
守恒检查以各网格自己的初始积分为基准；初始近似积分对非多项式涡的误差，
不能误计为时间推进中的守恒损失。完整动量漂移、密度与压力正性检查、
配置和误差元数据保存在 \texttt{metrics.json}。

\section{最终场对比}
\begin{figure}[H]\centering
\includegraphics[width=\linewidth]{density_grid_comparison.pdf}
\caption{$t=2,\ z=2$ 的恢复节点密度与误差。
上排密度共用色标；下排每个误差图采用各自对称色标，
标题给出范围，不能仅凭颜色判断不同网格误差大小。
显示使用原始网格三角面上的线性着色；范数直接在原始节点计算，不使用插值图。}
\end{figure}

\section{复现与后续判断}
新增运行配置位于 \path{cases/NCFV/} 目录，文件名分别为
\path{NCFV_iv10.json}、\path{NCFV_iv20.json} 和
\path{NCFV_iv80.json}。

时间步检查使用同目录的 \path{NCFV_iv20_halfdt.json} 和
\path{NCFV_iv40_halfdt.json}。从 \path{build/} 启动相应配置。
所有新运行的日志、首末节点表、VTK 场和 H5 重启文件保存在
\path{data/out/NCFV/convergence/}。
iv40 基准结果沿用 \path{data/out/NCFV/iv40_efficient/}。

从仓库根目录执行后处理脚本：
\begin{verbatim}
venv/bin/python cases/NCFV/analyze_vortex_convergence.py
xelatex -interaction=nonstopmode -halt-on-error \
  -output-directory=docs/reports \
  docs/reports/NCFV_convergence_report.tex
\end{verbatim}
本次验证回答的是无黏等熵涡的当前高效实现精度，而不是黏性项、限制器、
复杂物理边界或大规模 MPI 扩展性能的综合验证。若需要严谨建立稳定渐近三阶，
可继续细化网格，并在最细网格补做时间步敏感性与独立光滑制造解检查。
\end{document}
